\PassOptionsToPackage{colorlinks=true,urlcolor=navyblue}{hyperref}
\documentclass[11pt,a4paper,sans]{moderncv}
\moderncvstyle{classic}
\moderncvcolor{blue}

\usepackage[utf8]{inputenc}
\usepackage[top=1.3cm,  bottom=1.2cm, left=1.5cm, right=1.5cm]{geometry}
\usepackage{multicol}

\definecolor{navyblue}{RGB}{0, 0, 128}

\colorlet{bodyrulecolor}{blue}
\colorlet{sectioncolor}{blue}

\usepackage{enumitem}  % for item sizes
\setlist[itemize]{topsep=3pt, itemsep=2pt, parsep=0pt}   % for item sizes

\name{Javad}{Komijani}
\title{Research Scientist, Generative Models \& Deep Learning}
\address{Zurich, Switzerland}
\phone[mobile]{+41 76 796 9992}
\email{jkomijani@gmail.com}
\homepage{jkomijani.github.io}
\social[github]{jkomijani}
% \social[linkedin]{javad-komijani-543a31235}


\begin{document}
\makecvtitle


\section{Professional Summary}

Research scientist in generative modeling with a background in theoretical and computational physics.
Since 2022, developing diffusion, flow-matching, and normalizing-flow models that respect exact symmetries, for sampling quantum field theories on grids, together with open-source PyTorch software that supports distributed training.
Earlier work in precision lattice QCD and mathematical physics brings the mathematical depth to design models and make them reliable.

\vspace{0.5em} \textbf{Research impact:} 27 peer-reviewed publications (\href{https://jkomijani.github.io/docs/Publications.pdf}{full list}), h-index 20, 3426 citations (\href{https://scholar.google.com/citations?user=EBV_Sl0AAAAJ&hl=en}{Google Scholar}, Oct 2026).

%--------------------------------------------------
\section*{Research Highlights}

\begin{itemize}[leftmargin=*]
    \item \textbf{Parameter-efficient generative architectures} for physics simulations, where Monte Carlo slows down near critical points, building in exact symmetries:
    \begin{itemize}
    \item PSD flow: a Fourier-space layer inspired by effective field theory that reshapes correlations at all scales at once, in a model with only $\sim$3.4K parameters.
    \item SVD-based gauge-equivariant flows: 
    a 1288-parameter block on a $4^4$ SU(3) lattice achieves effective sample size comparable to a 288-layer spectral-flow baseline.
    \end{itemize}
    \item \textbf{Symmetry-aware diffusion models:}
    \begin{itemize}
    \item Equivariant diffusion: reverse-time SDEs from an exact-likelihood ODE to Langevin-corrected sampling; much shorter autocorrelation than Hybrid Monte Carlo (HMC).
    \item Diffusion on Lie groups: implicit score matching; the SDE family extended to Lie groups; compact gauge-equivariant networks and a multi-scale U-Net.
    \item Noise schedule on Lie groups: derived analytically; gives linear energy decay and few-step sampling without the hand-designed drift Euclidean models need.
    \end{itemize}
    \item \textbf{Machine learning combined with statistics}: symmetries split a sampling problem into independent pieces learned by a tiny flow, plus an exact statistical reweighting that removes the bias; validated against HMC.
    \item \textbf{Precision physics:} High-precision determinations of Standard Model parameters (quark masses, decay constants) from large-scale lattice simulations in international collaborations, using effective field theories; results in the particle-physics world averages (FLAG, PDG).
    \item \textbf{Mathematical physics and hyperasymptotics:} Resummation of divergent perturbative series, including the minimal renormalon-subtracted (MRS) scheme used for precise heavy-quark masses; methods for nonlinear eigenvalue problems.
\end{itemize}


\section*{Research \& Technical Experience}

\cventry{2021--2026}{Postdoctoral Researcher}{ETH Zurich}{}{}{
\begin{itemize}
\item Developed generative-model methods, open-source software, and distributed-training infrastructure.
\item Integrated normalizing-flow modules into the EU \href{https://www.intertwin.eu/article/thematic-module-normflow}{interTwin Digital Twin Engine}.
\item Co-supervised PhD, MSc, BSc, and intern projects in computational physics and scientific ML.
\end{itemize}
}

\cventry{2015--2021}{Postdoctoral Researcher}{TU Munich / Glasgow / Tehran}{}{}{
\begin{itemize}
\item Monte Carlo simulations and Bayesian parameter estimation for precision lattice-QCD observables.
\item Developed scalable lattice-QCD simulation pipelines on HPC systems in international collaborations.
\end{itemize}
}


\clearpage
%--------------------------------------------------
\section{Open-Source Scientific ML Software}

\cventry{\href{https://github.com/jkomijani/lattice_ml}{lattice\_ml}}{Composable PyTorch library for lattice field theory (v0.2.0)}{}{}{}{
\begin{itemize}
\item Adjoint-method ODE integrators for memory-efficient backpropagation of derivatives through long trajectories, compatible with automatic differentiation, for scalar and Lie-group-valued objects; integrators that stay on the group manifold; symplectic solvers; Hutchinson trace estimation.
\item Implicit score matching and diffusion on Lie groups with predictor--corrector sampling; flow matching and flow maps.
\item HMC for U(1) and SU($N$); AD-safe SVD and eigendecomposition; gauge-equivariant layers; prelink and holonomy tools; Lie-group coordinates and random-matrix sampling; distributed training.
\item Used in arXiv:2605.06134, 2605.17326, and 2610.09147 (\href{https://github.com/jkomijani/holonomy-2d-gauge-notebooks}{reproducible notebooks}).
\end{itemize}
}
\cventry{\href{https://github.com/jkomijani/normflow}{normflow}}{PyTorch package for normalizing flows in lattice field theory (v3.0.0, since 2021)}{}{}{}{
\begin{itemize}
\item Scalar and SU($N$) gauge theories; self-learning from the action (reverse KL with path gradients) or data-based training; asymptotically exact sampling via a Metropolis step.
\item EFT-inspired PSD and Fourier-space flows; coupling, rational-quadratic spline, and autoregressive layers.
\item Spectral and SVD-based gauge-equivariant flows; distributed training.
\end{itemize}
}


%--------------------------------------------------
\section{Technical Skills}

\cvitem{Programming}{Python, C++, Cython, Bash; Linux, Git}
\cvitem{ML/AI}{PyTorch, Diffusion Models, Flow Matching, Normalizing Flows, Score Matching, Equivariant Deep Learning, Efficient Sampling (HMC, Predictor--Corrector), Distributed Training}
\cvitem{Mathematics}{Differential Equations, Linear Algebra, Lie Groups, Asymptotic Analysis, Optimization}
\cvitem{Statistics}{Monte Carlo Methods, Bayesian Inference, Importance Sampling, Statistical Modeling}
\cvitem{Computing}{CUDA, MPI, HPC and Supercomputing Environments, NumPy, SciPy}
\cvitem{Languages}{English (fluent), German (basic, A2)}


%--------------------------------------------------
\section*{Mentoring \& Leadership}

\cventry{Mentoring}{Co-supervised PhD, MSc, BSc, and intern researchers in physics \& scientific ML:}{}{}{}{
\begin{itemize}
\item PhD (2): gauge-equivariant diffusion models; Fourier acceleration for SU($N$).
\item MSc (6): equivariant U-Net; autoregressive, continuous-time, and spline-based flows; simulated tempering.
\item BSc and intern: path-integral Monte Carlo; CNNs in arbitrary dimensions.
\end{itemize}
}

\cventry{Organization}{Co-organized workshops on physics \& machine learning:}{}{}{}{
\begin{itemize}
\item ZPW 2026 --- Machine Learning Techniques for Particle Physics.
\item QCD 2025 --- Symposium on QCD Factorization.
\item ML Lattice 2025 --- Machine learning approaches in Lattice QCD.
\item UT 2019 --- Lattice QCD workshop; TUM 2017 --- EFT and Lattice Field Theory Summer School.
\end{itemize}
}

\cvitem{Refereeing}{Referee for Physical Review D (13+ reviews since 2016).}

%--------------------------------------------------
\section{Education}

\cventry{2015}{PhD in Physics}{Washington University in St. Louis}{}{}{
\textbf{Thesis}: \textit{Topics in Lattice Gauge Theory and Theoretical Physics}\\
\textbf{Advisors}: Claude Bernard and Carl M. Bender
\begin{itemize}
  \item Developed EFT-based frameworks and Bayesian inference methods for lattice QCD.
  \item Performed precision determinations of hadronic observables in lattice QCD.
  \item Developed a novel definition of eigenvalue problems for nonlinear systems arising in physics.
\end{itemize}
}
\cventry{2009}{MSc in Electrical Engineering}{University of Tehran}{}{}{
\begin{itemize}
\item Focus: electromagnetic waves and fields, signal processing, statistical analysis.
\end{itemize}
}
\cventry{2006}{BSc in Electrical Engineering}{University of Tehran}{}{}{}

\end{document}